\section{Pentose Phosphate Pathway}
%%%%%%%%%%%%%%%%%%%%%%%%%%%%%%%%%%%%%%%%%%%%%%%%%%%%%%%%%%%%%%%

\subsection{Learning Objectives}
\begin{itemize}
    \item Describe the two phases of the pentose phosphate pathway
    \item Explain the production and importance of NADPH
    \item Understand the synthesis of ribose-5-phosphate for nucleotide biosynthesis
    \item Describe the regulation of the pathway
    \item Explain the clinical significance of G6PD deficiency
\end{itemize}

\subsection{Overview and Functions}

\begin{definition}
The \textbf{pentose phosphate pathway} (PPP), also called the \textbf{hexose monophosphate shunt}, is an alternative route for glucose-6-phosphate oxidation that produces:
\begin{enumerate}
    \item \textbf{NADPH}: Reducing power for biosynthesis and antioxidant defense
    \item \textbf{Ribose-5-phosphate}: Precursor for nucleotide synthesis
\end{enumerate}
\end{definition}

\textbf{Location:} Cytosol (same as glycolysis)

\textbf{Tissues with high PPP activity:}
\begin{itemize}
    \item \textbf{Liver}: Fatty acid and cholesterol synthesis (NADPH)
    \item \textbf{Adipose tissue}: Fatty acid synthesis (NADPH)
    \item \textbf{Adrenal cortex}: Steroid hormone synthesis (NADPH)
    \item \textbf{Red blood cells}: Glutathione reduction (NADPH)
    \item \textbf{Rapidly dividing cells}: Nucleotide synthesis (ribose-5-P)
\end{itemize}

\subsection{Oxidative Phase (Irreversible)}

The oxidative phase produces NADPH and converts glucose-6-phosphate to ribulose-5-phosphate.

\textbf{Reaction 1: Glucose-6-phosphate Dehydrogenase (G6PD)}

\[
\text{Glucose-6-P} + \text{NADP}^+ \xrightarrow{\text{G6PD}} \text{6-Phosphogluconolactone} + \text{NADPH} + \text{H}^+
\]

\begin{itemize}
    \item \textbf{Rate-limiting step} of the PPP
    \item \textbf{First NADPH produced}
    \item Irreversible; commits G6P to the PPP
    \item \textbf{Regulated by NADP$^+$/NADPH ratio}
\end{itemize}

\textbf{Reaction 2: Lactonase}

\[
\text{6-Phosphogluconolactone} + \text{H}_2\text{O} \rightarrow \text{6-Phosphogluconate}
\]

\begin{itemize}
    \item Hydrolysis of the lactone ring
    \item Spontaneous but enzyme-catalyzed for speed
\end{itemize}

\textbf{Reaction 3: 6-Phosphogluconate Dehydrogenase}

\[
\text{6-Phosphogluconate} + \text{NADP}^+ \rightarrow \text{Ribulose-5-P} + \text{CO}_2 + \text{NADPH}
\]

\begin{itemize}
    \item Oxidative decarboxylation
    \item \textbf{Second NADPH produced}
    \item Converts 6C sugar to 5C sugar (releases CO$_2$)
\end{itemize}

\begin{tcolorbox}[colback=blue!5, colframe=blue!70!black, title=\textbf{Oxidative Phase Summary}]
\textbf{Input:} Glucose-6-phosphate + 2 NADP$^+$ + H$_2$O

\textbf{Output:} Ribulose-5-phosphate + 2 NADPH + 2 H$^+$ + CO$_2$

\textbf{Key points:}
\begin{itemize}
    \item Irreversible (committed pathway
)
    \item Generates reducing power (NADPH)
    \item Loss of one carbon as CO$_2$
\end{itemize}
\end{tcolorbox}

\subsection{Non-Oxidative Phase (Reversible)}

The non-oxidative phase interconverts sugars of different carbon lengths through a series of reversible reactions.

\textbf{Key Enzymes:}

\begin{enumerate}
    \item \textbf{Phosphopentose Isomerase}
    \[
    \text{Ribulose-5-P} \xleftrightarrow{} \text{Ribose-5-P}
    \]
    \begin{itemize}
        \item Ketose to aldose isomerization
        \item Produces ribose-5-P for nucleotide synthesis
    \end{itemize}
    
    \item \textbf{Phosphopentose Epimerase}
    \[
    \text{Ribulose-5-P} \xleftrightarrow{} \text{Xylulose-5-P}
    \]
    \begin{itemize}
        \item Epimerization at C3
        \item Generates substrate for transketolase
    \end{itemize}
    
    \item \textbf{Transketolase} (requires TPP)
    \begin{itemize}
        \item Transfers 2-carbon units
        \item Thiamine pyrophosphate (vitamin B$_1$) is cofactor
    \end{itemize}
    
    \item \textbf{Transaldolase}
    \begin{itemize}
        \item Transfers 3-carbon units
        \item No cofactor required
    \end{itemize}
\end{enumerate}

\textbf{Carbon Shuffling Summary:}

\begin{center}
\begin{tikzpicture}[node distance=1.5cm]
    \node (r5p1) {Ribose-5-P (5C)};
    \node[right=2cm of r5p1] (x5p1) {Xylulose-5-P (5C)};
    \node[below=1cm of r5p1, xshift=1.5cm] (middle1) {};
    \node[below=0.5cm of middle1] (s7p) {Sedoheptulose-7-P (7C)};
    \node[right=2cm of s7p] (g3p1) {G3P (3C)};
    \node[below=1cm of s7p] (e4p) {Erythrose-4-P (4C)};
    \node[right=2cm of e4p] (f6p1) {Fructose-6-P (6C)};
    \node[below=1cm of e4p, xshift=1.5cm] (middle2) {};
    \node[below=0.5cm of middle2] (f6p2) {Fructose-6-P (6C)};
    \node[right=2cm of f6p2] (g3p2) {G3P (3C)};
    
    \draw[->, thick] (r5p1) -- node[above] {\small TK} (x5p1);
    \draw[->, thick] (x5p1) -- (s7p);
    \draw[->, thick] (r5p1) -- (g3p1);
    \draw[->, thick] (s7p) -- node[left] {\small TA} (e4p);
    \draw[->, thick] (g3p1) -- (f6p1);
    \draw[->, thick] (e4p) -- node[above] {\small TK} (f6p2);
\end{tikzpicture}
\end{center}

\textbf{Net reaction of non-oxidative phase:}
\[
3 \text{ Ribulose-5-P (5C)} \xleftrightarrow{} 2 \text{ Fructose-6-P (6C)} + 1 \text{ G3P (3C)}
\]

\begin{itemize}
    \item 15 carbons in = 15 carbons out (12 + 3)
    \item F6P and G3P can enter glycolysis or gluconeogenesis
    \item Reversibility allows pathway to meet cellular needs
\end{itemize}

\subsection{Modes of the Pentose Phosphate Pathway}

\begin{theorem}
The PPP operates in different modes depending on cellular needs:

\textbf{Mode 1: Need NADPH and Ribose-5-P (balanced)}
\begin{itemize}
    \item Oxidative phase runs
    \item Ribulose-5-P $\rightarrow$ Ribose-5-P (isomerase only)
    \item Example: Rapidly dividing cells
\end{itemize}

\textbf{Mode 2: Need NADPH only (excess ribose)}
\begin{itemize}
    \item Oxidative phase runs maximally
    \item Non-oxidative phase converts R5P $\rightarrow$ F6P + G3P
    \item F6P and G3P enter glycolysis
    \item Example: Adipose tissue during lipogenesis
\end{itemize}

\textbf{Mode 3: Need Ribose-5-P only (no NADPH need)}
\begin{itemize}
    \item Oxidative phase bypassed
    \item Non-oxidative phase runs in reverse
    \item F6P + G3P $\rightarrow$ Ribose-5-P
    \item Example: Rapidly dividing cells with adequate NADPH
\end{itemize}

\textbf{Mode 4: Need NADPH and ATP}
\begin{itemize}
    \item Oxidative phase produces NADPH
    \item R5P converted to F6P and G3P
    \item These enter glycolysis for ATP production
    \item Example: RBCs (need NADPH for glutathione, ATP for ion pumps)
\end{itemize}
\end{theorem}

\subsection{Regulation of the Pentose Phosphate Pathway}

\begin{itemize}
    \item \textbf{Primary control:} G6PD (rate-limiting enzyme)
    \item \textbf{Activated by:} High NADP$^+$/NADPH ratio (need for NADPH)
    \item \textbf{Inhibited by:} High NADPH (product inhibition)
\end{itemize}

\begin{tcolorbox}[colback=red!5, colframe=red!70!black, title=\textbf{Clinical Correlation: G6PD Deficiency}]
\textbf{G6PD deficiency} is the most common enzyme deficiency worldwide, affecting over 400 million people.

\textbf{Pathophysiology:}
\begin{itemize}
    \item Decreased NADPH production in RBCs
    \item Impaired glutathione reduction (GSH maintenance)
    \item RBCs vulnerable to oxidative stress
    \item Hemolysis upon exposure to oxidants
\end{itemize}

\textbf{Triggers of hemolytic crisis:}
\begin{itemize}
    \item Fava beans (favism)
    \item Drugs: Primaquine, sulfonamides, dapsone, nitrofurantoin
    \item Infections (increased oxidative stress)
    \item Diabetic ketoacidosis
\end{itemize}

\textbf{Clinical features:}
\begin{itemize}
    \item Acute hemolytic anemia
    \item Jaundice, dark urine (hemoglobinuria)
    \item ``Bite cells'' and Heinz bodies on blood smear
\end{itemize}

\textbf{Inheritance:} X-linked recessive (males predominantly affected)

\textbf{Protective effect:} Confers resistance to \textit{Plasmodium falciparum} malaria
\end{tcolorbox}

\subsection{Functions of NADPH}

\begin{enumerate}
    \item \textbf{Reductive biosynthesis:}
    \begin{itemize}
        \item Fatty acid synthesis
        \item Cholesterol synthesis
        \item Steroid hormone synthesis
        \item Neurotransmitter synthesis
    \end{itemize}
    
    \item \textbf{Antioxidant defense:}
    \begin{itemize}
        \item Regeneration of reduced glutathione (GSH)
        \item Thioredoxin reduction
        \item Protection against reactive oxygen species
    \end{itemize}
    
    \item \textbf{Detoxification:}
    \begin{itemize}
        \item Cytochrome P450 reactions
        \item Drug metabolism in liver
    \end{itemize}
    
    \item \textbf{Immune function:}
    \begin{itemize}
        \item NADPH oxidase in phagocytes
        \item Respiratory burst (superoxide production)
    \end{itemize}
\end{enumerate}

\begin{example}
\textbf{Glutathione System in RBCs}

Red blood cells lack mitochondria and rely entirely on the PPP for NADPH production.

\textbf{The glutathione cycle:}
\begin{enumerate}
    \item Glutathione peroxidase uses GSH to reduce H$_2$O$_2$:
    \[
    2\text{GSH} + \text{H}_2\text{O}_2 \xrightarrow{\text{GPx}} \text{GSSG} + 2\text{H}_2\text{O}
    \]
    
    \item Glutathione reductase regenerates GSH using NADPH:
    \[
    \text{GSSG} + \text{NADPH} + \text{H}^+ \xrightarrow{\text{GR}} 2\text{GSH} + \text{NADP}^+
    \]
\end{enumerate}

Without adequate NADPH (as in G6PD deficiency), GSH cannot be regenerated, and oxidative damage accumulates.
\end{example}

\section{Glycogen Metabolism}

\subsection{Overview of Glycogen}

\begin{definition}
\textbf{Glycogen} is a highly branched polymer of glucose that serves as the primary storage form of glucose in animals. It is stored mainly in:
\begin{itemize}
    \item \textbf{Liver} ($\sim$100g): Maintains blood glucose levels
    \item \textbf{Skeletal muscle} ($\sim$400g): Fuel for muscle contraction
\end{itemize}
\end{definition}

\textbf{Structure of Glycogen:}
\begin{itemize}
    \item Linear chains: $\alpha$(1$\rightarrow$4) glycosidic bonds
    \item Branch points: $\alpha$(1$\rightarrow$6) glycosidic bonds (every 8-12 residues)
    \item Highly branched structure allows rapid mobilization
    \item Each glycogen granule contains up to 55,000 glucose residues
    \item \textbf{Glycogenin} protein at the core (primer for synthesis)
\end{itemize}

\begin{analogy}
Glycogen is like a tree with many branches. The trunk represents the core protein (glycogenin), and the branches ($\alpha$-1,6 linkages) allow many ``leaves'' (glucose residues) to be added or removed simultaneously, enabling rapid storage or release of glucose.
\end{analogy}

\subsection{Glycogenesis (Glycogen Synthesis)}

\begin{definition}
\textbf{Glycogenesis} is the process of synthesizing glycogen from glucose. It occurs primarily in liver and muscle when glucose is abundant (fed state).
\end{definition}

\textbf{Steps of Glycogenesis:}

\textbf{Step 1: Glucose Phosphorylation}
\[
\text{Glucose} + \text{ATP} \xrightarrow{\text{Hexokinase/Glucokinase}} \text{Glucose-6-P} + \text{ADP}
\]

\textbf{Step 2: Isomerization}
\[
\text{Glucose-6-P} \xleftrightarrow{\text{Phosphoglucomutase}} \text{Glucose-1-P}
\]

\textbf{Step 3: Activation (UDP-glucose formation)}
\[
\text{Glucose-1-P} + \text{UTP} \xrightarrow{\text{UDP-glucose pyrophosphorylase}} \text{UDP-glucose} + \text{PP}_i
\]
\begin{itemize}
    \item UDP-glucose is the ``activated'' form of glucose
    \item Pyrophosphate (PP$_i$) hydrolysis makes reaction irreversible
    \item PP$_i$ $\rightarrow$ 2 P$_i$ (by pyrophosphatase)
\end{itemize}

\textbf{Step 4: Chain Elongation}
\[
\text{UDP-glucose} + \text{Glycogen}_{(n)} \xrightarrow{\text{Glycogen synthase}} \text{Glycogen}_{(n+1)} + \text{UDP}
\]
\begin{itemize}
    \item \textbf{Glycogen synthase}: Key regulatory enzyme
    \item Adds glucose to non-reducing end via $\alpha$(1$\rightarrow$4) bond
    \item Requires pre-existing primer (at least 4 glucose residues)
\end{itemize}

\textbf{Step 5: Branch Formation}
\[
\text{Linear chain} \xrightarrow{\text{Branching enzyme}} \text{Branched chain}
\]
\begin{itemize}
    \item \textbf{Branching enzyme} (glycosyl 4:6 transferase)
    \item Transfers 6-7 glucose block from end of chain
    \item Creates new $\alpha$(1$\rightarrow$6) branch point
    \item Branch must be at least 4 residues from existing branch
\end{itemize}

\textbf{Primer Synthesis (Glycogenin):}
\begin{itemize}
    \item Glycogenin is a protein that self-glucosylates
    \item Attaches first glucose to tyrosine residue
    \item Extends chain to $\sim$8 glucose residues
    \item Then glycogen synthase takes over
\end{itemize}

\begin{tcolorbox}[colback=green!5, colframe=green!70!black, title=\textbf{Energy Cost of Glycogenesis}]
\textbf{Per glucose added to glycogen:}
\begin{itemize}
    \item 1 ATP (hexokinase)
    \item 1 UTP (UDP-glucose pyrophosphorylase)
\end{itemize}
\textbf{Total: 2 ATP equivalents per glucose stored}
\end{tcolorbox}

\subsection{Glycogenolysis (Glycogen Breakdown)}

\begin{definition}
\textbf{Glycogenolysis} is the process of breaking down glycogen to release glucose (liver) or glucose-6-phosphate (muscle) for energy production.
\end{definition}

\textbf{Steps of Glycogenolysis:}

\textbf{Step 1: Chain Shortening}
\[
\text{Glycogen}_{(n)} + \text{P}_i \xrightarrow{\text{Glycogen phosphorylase}} \text{Glycogen}_{(n-1)} + \text{Glucose-1-P}
\]
\begin{itemize}
    \item \textbf{Glycogen phosphorylase}: Key regulatory enzyme
    \item Cleaves $\alpha$(1$\rightarrow$4) bonds by phosphorolysis
    \item Releases glucose-1-phosphate (not free glucose)
    \item Stops 4 residues from branch point
    \item Requires pyridoxal phosphate (PLP, vitamin B$_6$) cofactor
\end{itemize}

\textbf{Step 2: Branch Removal (Debranching Enzyme)}

The debranching enzyme has two activities:

\begin{enumerate}
    \item \textbf{Transferase activity:}
    \begin{itemize}
        \item Moves 3 of the 4 remaining glucose residues to another chain
        \item Creates new $\alpha$(1$\rightarrow$4) linkage
    \end{itemize}
    
    \item \textbf{$\alpha$-1,6-Glucosidase activity:}
    \begin{itemize}
        \item Hydrolyzes the $\alpha$(1$\rightarrow$6) bond
        \item Releases single free glucose molecule
        \item This is the only free glucose released in glycogenolysis
    \end{itemize}
\end{enumerate}

\textbf{Step 3: Conversion to Glucose-6-P}
\[
\text{Glucose-1-P} \xleftrightarrow{\text{Phosphoglucomutase}} \text{Glucose-6-P}
\]

\textbf{Step 4: Fate of Glucose-6-P}

\begin{itemize}
    \item \textbf{In muscle:} G6P enters glycolysis (muscle lacks G6Pase)
    \item \textbf{In liver:} G6P $\rightarrow$ Glucose (via glucose-6-phosphatase)
    \[
    \text{Glucose-6-P} + \
text{H}_2\text{O} \xrightarrow{\text{Glucose-6-phosphatase}} \text{Glucose} + \text{P}_i
    \]
    \begin{itemize}
        \item Located in ER membrane
        \item Allows liver to release glucose into blood
    \end{itemize}
\end{itemize}

\begin{tcolorbox}[colback=yellow!10, colframe=yellow!70!black, title=\textbf{Glycogenolysis Products}]
\textbf{Approximately 90\%:} Glucose-1-phosphate (from $\alpha$-1,4 cleavage)

\textbf{Approximately 10\%:} Free glucose (from $\alpha$-1,6 branch points)
\end{tcolorbox}

\subsection{Regulation of Glycogen Metabolism}

Glycogen metabolism is regulated by reciprocal control of glycogen synthase and glycogen phosphorylase through:
\begin{itemize}
    \item Hormonal signals (insulin, glucagon, epinephrine)
    \item Allosteric effectors
    \item Covalent modification (phosphorylation)
\end{itemize}

\textbf{Hormonal Regulation:}

\begin{center}
\begin{tabular}{lccc}
\toprule
\textbf{Hormone} & \textbf{Target Tissue} & \textbf{Glycogen Synthase} & \textbf{Glycogen Phosphorylase} \\
\midrule
Insulin & Liver, Muscle & Activates & Inhibits \\
Glucagon & Liver only & Inhibits & Activates \\
Epinephrine & Liver, Muscle & Inhibits & Activates \\
\bottomrule
\end{tabular}
\end{center}

\textbf{Phosphorylation States:}

\begin{itemize}
    \item \textbf{Glycogen synthase:}
    \begin{itemize}
        \item Dephosphorylated = Active (GSa)
        \item Phosphorylated = Inactive (GSb)
    \end{itemize}
    
    \item \textbf{Glycogen phosphorylase:}
    \begin{itemize}
        \item Dephosphorylated = Inactive (GPb)
        \item Phosphorylated = Active (GPa)
    \end{itemize}
\end{itemize}

\begin{tcolorbox}[colback=red!5, colframe=red!70!black, title=\textbf{Clinical Correlation: Glycogen Storage Diseases}]
\textbf{Von Gierke Disease (Type I):}
\begin{itemize}
    \item Deficiency: Glucose-6-phosphatase
    \item Features: Severe fasting hypoglycemia, hepatomegaly, lactic acidosis
    \item Liver cannot release free glucose
\end{itemize}

\textbf{Pompe Disease (Type II):}
\begin{itemize}
    \item Deficiency: Lysosomal $\alpha$-1,4-glucosidase (acid maltase)
    \item Features: Cardiomegaly, muscle weakness, early death
    \item Glycogen accumulates in lysosomes
\end{itemize}

\textbf{McArdle Disease (Type V):}
\begin{itemize}
    \item Deficiency: Muscle glycogen phosphorylase
    \item Features: Exercise intolerance, muscle cramps, myoglobinuria
    \item Cannot mobilize muscle glycogen during exercise
\end{itemize}

\textbf{Cori Disease (Type III):}
\begin{itemize}
    \item Deficiency: Debranching enzyme
    \item Features: Milder hypoglycemia, limit dextrin accumulation
    \item Abnormal glycogen structure with short outer branches
\end{itemize}
\end{tcolorbox}

